\section{评估与基准}

\subsection{多层评估框架}

讲者把评估分成三层：离线基准（benchmarks）、模拟部署（sandbox）、实际回放（replay evaluation）。

\begin{itemize}
    \item \textbf{离线 benchmark}：D4RL 类型的任务，关注 normalized score。
    \item \textbf{模拟部署}：用 simulator 跑 policy with hold-out reward to detect inflated Q-values。
    \item \textbf{回放评估}：把 policy 产生的 trajectory 反写入日志，检查 reward/constraint 是否满足。
\end{itemize}

\begin{importantbox}{评估越多越好，但须留出 human-in-the-loop}
自动 benchmark 只能检验 capacity；要理解 failure modes，还需要串联 human audit：观察 policy 在关键情景（edge case）下的行为，补充 qualitative assessment。
\end{importantbox}

\begin{figure}[H]
\centering
\includegraphics[page=3,width=0.9\textwidth]{07_cs224r_offline_rl_2025.pdf}
\caption{Lecture slides 描述的三层评估管线\protect\footnotemark}
\end{figure}
\footnotetext{Slides 第3页，高亮 offline RL 的评估阶梯。}

\subsection{评估中的对比表}

\begin{table}[H]
\centering
\begin{tabular}{p{0.2\textwidth} p{0.26\textwidth} p{0.26\textwidth}}
\toprule
评估阶层 & 主要目的 & 常用工具 \\
\midrule
离线 benchmark & 快速对齐整体能力 & D4RL, OpenAI 的 offline suites \\
Replay evaluation & 观察 policy 在 replay buffer 中的 reward & replay simulator, human audit \\
Sandbox deployment & 用 simulator 到达准生产环境 & per-domain simulator + twin models \\
\bottomrule
\end{tabular}
\caption{评估层级的对照表}
\end{table}

\subsection{本章小结}

评估应该包括 benchmark、sandbox 尝试和回放验证三层，每层都可捕捉不同 failure 模式。人类审查仍不可或缺。

